\documentclass[a4paper, serif, 10pt]{moderncv}

%% moderncv
% style and color
\moderncvstyle{banking}
\moderncvcolor{black}

%% page layout/margins
\usepackage[top=2.5cm, bottom=2.5cm, left=1.5cm, right=1.5cm]{geometry}

%% Babel config for German language
% ref:
% - https://latex3.github.io/babel/guides/locale-german.html
\usepackage[ngerman]{babel}

%% fontspec
\usepackage{fontspec}

\IfFontExistsTF{Linux Libertine}{
  \setmainfont[Ligatures={TeX, Common}, Numbers=OldStyle]{Linux Libertine}
}{}
\IfFontExistsTF{Linux Libertine O}{
  \setmainfont[Ligatures={TeX, Common}, Numbers=OldStyle]{Linux Libertine O}
}{}

%% CTeX
% CJK support, used to show CJK characters in the resume
%
% - fontset=none: disable builtin fontset but instead set the CJK font manually
% - heading=false: disable ctex heading
% - punct=kaiming: use kaiming punctuations styles for CJK
% - scheme=plain: use plain scheme, do not override \normalsize font size
% - space=auto: space settings for CJK characters
%
% ref:
% - http://ctan.mirrorcatalogs.com/language/chinese/ctex/ctex.pdf
\usepackage[UTF8, heading=false, punct=kaiming, scheme=plain, space=auto]{ctex}

\IfFontExistsTF{Noto Serif CJK SC}{
  \setCJKmainfont{Noto Serif CJK SC}
}{}
\IfFontExistsTF{Noto Sans CJK SC}{
  \setCJKsansfont{Noto Sans CJK SC}
}{}

%% line spacing
\usepackage{setspace}
\setstretch{1.125}

%% URL styling - use normal text instead of monospace
\urlstyle{same}

%% Patch \httplink and \httpslink to handle full URLs
% The original moderncv macros always prepend http:// or https:// to URLs.
% This causes issues when the URL already has a protocol (e.g., https://example.com)
% resulting in malformed URLs like https://https://example.com.
% This patch checks if the URL already contains "://" and uses it directly if so.
% Note: We use \str_set:Nx (x-expansion) to expand macros like \@homepage first.
\ExplSyntaxOn
\str_new:N \g_yamlresume_url_str

\RenewDocumentCommand{\httpslink}{O{}m}{
  \str_set:Nx \g_yamlresume_url_str {#2}
  \str_if_in:NnTF \g_yamlresume_url_str {://}
    { \url{#2} }
    { \url{https://#2} }
}

\RenewDocumentCommand{\httplink}{O{}m}{
  \str_set:Nx \g_yamlresume_url_str {#2}
  \str_if_in:NnTF \g_yamlresume_url_str {://}
    { \url{#2} }
    { \url{http://#2} }
}
\ExplSyntaxOff

\name{Yunlong Li}{}
\title{KI- \& IoT-Ingenieur}
\phone[mobile]{+41 79 539 72 86}
\email{ddlong68@outlook.com}
\homepage{https://www.linkedin.com/in/yunlong-li-6891a72a3/}

\address{Zürich, Schweiz}{}{}

\begin{document}

\maketitle

\section{Basisinformationen}

\cvline{}{Ingenieur an der Schnittstelle von KI, IoT und technischen Systemen — vom
Algorithmusentwurf bis zur Produktlieferung. Leitung funktionsübergreifender
Teams, Koordination mit Industriepartnern und Umsetzung von Projekten vom
Forschungskonzept bis zum produktiven Einsatz. Verbindet fundierte technische
Kompetenz in ML/KI-Modellierung, eingebetteten Plattformen und datengetriebener
Steuerung mit praktischer Erfahrung in Teamaufbau, Stakeholder-Kommunikation
und durchgängigem Projektmanagement.}

\section{Ausbildung}

\cventry{Sept. 2022–Juni 2026}
        {Master of art, Computer Linguistics \& Neuroinformatics}
        {Universität Zürich (UZH)}
        {}
        {}
        {}

\section{Berufserfahrung}

\cventry{März 2025–Jan. 2026}
        {Gründer \& Entwickler – SEO-Automatisierungssystem}
        {Selbstständiges Startup}
        {}
        {}
        {\begin{itemize}
\item Gesamten Produktlebenszyklus verantwortet — von Marktanalyse und Systemarchitektur bis Entwicklung und Deployment.
\item KI-basiertes SEO-Automatisierungssystem mit intelligenten Agenten entworfen und implementiert, unter Einsatz von Vibe Coding für schnelles Prototyping.
\item Keyword-Analyse, Content-Optimierung und Performance-Monitoring mit Echtzeit-Analytik automatisiert.
\item Produkt-Roadmap, Nutzerfeedback-Integration und iterative Release-Zyklen eigenständig gesteuert.
\end{itemize}
\textbf{Schlüsselwörter}: KI-Agenten, SEO, Vibe Coding, Produktmanagement}

\cventry{Jan. 2025–Jan. 2026}
        {Gründer \& Produktentwickler – KI-gestützte Smart Cat Door \& IoT-System}
        {Selbstständiges Startup}
        {}
        {}
        {\begin{itemize}
\item End-to-End-Produktentwicklung geleitet — vom Konzeptentwurf über Hardware-Prototyping bis zum Cloud-Deployment.
\item Hardware (Embedded-Firmware), Software (Cloud-Backend) und Benutzeroberfläche koordiniert, mit Vibe Coding zur Beschleunigung der Full-Stack-Implementierung.
\item Embedded-KI für Tiererkennung und Verhaltenserkennung mit Micropython implementiert.
\item Plattformübergreifende Mobile-App mit Expo (React Native) für Haustierüberwachung und Türsteuerung entwickelt.
\item Gesamte Lieferkette aufgebaut — von Lieferantensuche und Komponentenbeschaffung bis zur Definition der Produktions-SOP.
\item Praxiswissen in CE-, RED- und Netzwerkzertifizierungen für den EU-Markteintritt erworben.
\item China-EU-Logistikkoordination inkl. Versand, Zoll und Lagerhaltung gesteuert.
\item Lieferantenkommunikation, Komponentenbeschaffung und iteratives Prototypen-Testing verantwortet.
\end{itemize}
\textbf{Schlüsselwörter}: IoT, Embedded-KI, Lieferkette, CE/RED-Zertifizierung, Produktentwicklung}

\cventry{Sept. 2023–Jan. 2024}
        {Forschungspraktikant Neurowissenschaften – In-Vivo-Neuralaufzeichnung}
        {HKUST (Guangzhou)}
        {}
        {}
        {\begin{itemize}
\item Mikrochirurgie (Kraniotomie) und Implantation von Neuralsonden für In-vivo-Aufzeichnungen durchgeführt.
\item Mit interdisziplinären Forschungsteams aus Neurowissenschaften und Ingenieurwesen zusammengearbeitet.
\end{itemize}
\textbf{Schlüsselwörter}: Neurowissenschaften, Neuralaufzeichnung, Interdisziplinäre Forschung}

\cventry{Sept. 2020–Juli 2022}
        {Praktikant HLK-Systementwicklung}
        {Guangzhou Design Institute}
        {}
        {}
        {\begin{itemize}
\item Projektteam aufgebaut und geleitet, Forschungsrichtung, Aufgabenverteilung und Meilensteinplanung definiert.
\item Algorithmen für multimodale Datenfusion entworfen und SAAS-Plattform-Architektur gestaltet.
\item Hauptansprechpartner zwischen Universitäts-Forschungsteam und Industriepartner — Auftragsannahme, Vor-Ort-Deployment und Projektabnahme gesteuert.
\item Projektpräsentationen und Roadshow-Pitches vor Stakeholdern und Investoren gehalten.
\item COP von 7,0 und 36 \% Energieeinsparung durch ML- \& IoT-getriebene Optimierung erreicht.
\end{itemize}
\textbf{Schlüsselwörter}: Teamführung, HLK, Stakeholder-Management, Projektlieferung}

\section{Sprachen}

\cvline{Mandarin}{Muttersprache}
\cvline{Kantonesisch}{Muttersprache}
\cvline{Englisch}{Verhandlungssicher}
\cvline{Deutsch}{Grundkenntnisse}

\section{Fähigkeiten}

\cvline{KI \& Maschinelles Lernen}{Experte \hfill \textbf{Schlüsselwörter}: Machine Learning, Deep Learning, Intelligente Agenten, RAG, Optimierungsalgorithmen, PyTorch}
\cvline{IoT, Embedded \& Engineering}{Erfahren \hfill \textbf{Schlüsselwörter}: Micropython, MQTT, Embedded-KI, Sensornetzwerke, HLK / PID / DDC, Gerät-Cloud-Kommunikation}
\cvline{Software \& Infrastruktur}{Erfahren \hfill \textbf{Schlüsselwörter}: Python, C/C++, JavaScript / TypeScript, Cloudflare (D1, R2, Worker, Network), Oracle Cloud, Docker / Nginx, N8N}
\cvline{Produkt \& Führung}{Erfahren \hfill \textbf{Schlüsselwörter}: Webentwicklung, Teamaufbau, Projektmanagement, Lieferkette / SOP, Stakeholder-Kommunikation}

\section{Veröffentlichungen}

\cventry{Feb. 2025}
        {DMPA: Model Poisoning Attacks on Decentralized Federated Learning for Model Differences}
        {arXiv}
        {\url{https://arxiv.org/abs/2502.04771}}
        {}
        {Innovative Model-Poisoning-Angriffsstrategie für dezentrale föderierte Lernsysteme vorgeschlagen.}

\cventry{Jan. 2022}
        {Real-time construction of thermal model based on multimodal scene data}
        {Frontiers in Energy Research (JCR Q2)}
        {}
        {}
        {Echtzeit-Thermomodellkonstruktion mittels multimodaler Szenendatenfusion entwickelt.}

\cventry{Jan. 2021}
        {Energy usage prediction based on multi-system data for public buildings using machine learning methods}
        {IEEE CompAuto 2021}
        {}
        {}
        {Machine-Learning-Methoden zur Energieverbrauchsvorhersage in öffentlichen Gebäuden angewendet.}

\cventry{Jan. 2020}
        {Online transaction detection method using CatBoost model}
        {IEEE CISCE 2020}
        {}
        {}
        {Online-Transaktionserkennung mittels CatBoost-Modell entwickelt.}

\section{Projekte}

\cventry{Sept. 2019–Jan. 2020}
        {Industrielles Bildverarbeitungssystem für flexible Fertigungs-Schraubmaschinen.}
        {Automatische Schraubmaschine – Bildverarbeitungssystem}
        {}
        {}
        {\begin{itemize}
\item Bildverarbeitungsmodul als alleiniger Verantwortlicher geleitet, Koordination mit Mechanik- und Embedded-Teams für Systemintegration.
\item Schraubenlochkoordinaten mittels 3D-Rekonstruktion extrahiert und an den Steuerrechner für automatisiertes Verschrauben übermittelt.
\item Limitation der 2D-Bildverarbeitung bei unterschiedlichen Positionen und Tiefen der Schraubenlöcher gelöst.
\item 3D-Modelle mittels Deep Learning und Punktwolkenverarbeitung erstellt; runde Schraubenlöcher mit Halcon identifiziert.
\item Erkennungsgeschwindigkeit von 130ms auf 30ms verbessert, direkter Durchsatzanstieg der Produktionslinie ermöglicht.
\end{itemize}
\textbf{Schlüsselwörter}: 3D-Rekonstruktion, Punktwolke, Halcon, Teamübergreifende Integration}

\cventry{Jan. 2020–Juli 2022}
        {Auf multimodaler Datenfusion basierendes HLK-Energiesparsystem.}
        {Energiesparendes Klimaanlagen-Steuerungssystem}
        {}
        {}
        {\begin{itemize}
\item Team von Grund auf aufgebaut — Mitglieder rekrutiert, Forschungsrichtung definiert und meilensteingesteuerte Projektpläne erstellt.
\item Algorithmus für multimodale Datenfusion entworfen und SAAS-Plattform für Echtzeit-Monitoring und -Steuerung architektiert.
\item Engineering-Lieferung durchgängig koordiniert: Unternehmenskommunikation, Auftragsannahme, Vor-Ort-Deployment und Abnahme.
\item IoT-Plattform und Sensordatenbank für Schulgebäude aufgebaut; multimodale Baseline-Kalibrierung durchgeführt.
\item Thermisches Komfort-Temperaturfeld-Analyseframework vorgeschlagen, 81 \% Genauigkeit und 36 \% Energieeinsparung erreicht.
\item Roadshow-Präsentationen und Projekt-Pitches vor Stakeholdern, Investoren und Jurymitgliedern gehalten.
\item Technologie im Produktivbetrieb eingesetzt und aktuell in Nutzung.
\end{itemize}
\textbf{Schlüsselwörter}: Teamführung, Multimodale Fusion, SAAS-Architektur, Stakeholder-Kommunikation, Projektlieferung}

\cventry{Jan. 2024–Feb. 2025}
        {Forschung zu Model-Poisoning-Angriffsstrategien für dezentrales föderiertes Lernen an der UZH.}
        {DMPA: Model-Poisoning-Angriffe auf dezentrales föderiertes Lernen}
        {}
        {}
        {\begin{itemize}
\item Paper mit internationalem Forschungsteam co-verfasst (arXiv:2502.04771).
\item Angriffsstrategie für dezentrales föderiertes Lernen durch Berechnung differenzieller Merkmale über adversariale Client-Modelle vorgeschlagen.
\item Experimente über mehrere Datensätze und Angriffsszenarien mit Co-Autoren koordiniert.
\item Methode übertrifft bestehende Model-Poisoning-Angriffe in dezentralen FL-Szenarien nachweislich.
\end{itemize}
\textbf{Schlüsselwörter}: Föderiertes Lernen, Model Poisoning, Forschungskooperation, Sicherheit}

\cventry{Juni 2025–Juni 2026}
        {Deep-Learning-basierte isolierte Gebärdenspracherkennung an der UZH.}
        {Gebärdenspracherkennung (ISLR) Forschung}
        {}
        {}
        {\begin{itemize}
\item Eigenständige Deep-Learning-Forschung zur isolierten Gebärdenspracherkennung (ISLR) durchgeführt.
\item Video-to-Pose-Pipeline zur Extraktion skelettaler Pose-Repräsentationen aus Gebärdensprachvideos untersucht.
\item Pose-to-Gloss-Mapping zur Übersetzung von Pose-Sequenzen in Gebärdensprach-Glossen erforscht.
\item Umfassende Forschungszusammenfassung mit Vergleich der Ergebnisse beider Richtungen erstellt.
\end{itemize}
\textbf{Schlüsselwörter}: Deep Learning, Gebärdensprache, Video-to-Pose, Pose-to-Gloss}

\cventry{März 2025–Jan. 2026}
        {KI-gesteuerte SEO-Automatisierungsplattform mit intelligenten Agenten.}
        {SEO-Automatisierungssystem}
        {}
        {}
        {\begin{itemize}
\item Systemarchitektur eigenständig entworfen, Produktanforderungen definiert und den gesamten Entwicklungslebenszyklus gesteuert.
\item KI-basiertes SEO-Automatisierungssystem mit intelligenten Agenten für Keyword-Analyse, Content-Optimierung und Performance-Monitoring aufgebaut.
\item Echtzeit-Analytik-Dashboards für datengetriebene Entscheidungsfindung integriert.
\item Nutzerakquise, Feedback-Sammlung und iterative Produktverbesserung verantwortet.
\end{itemize}
\textbf{Schlüsselwörter}: KI-Agenten, SEO, Produktverantwortung, Full-Stack}

\cventry{Jan. 2025–Jan. 2026}
        {KI-gestützte IoT-Plattform für intelligente Haustier-Zugangskontrolle.}
        {Smart Cat Door – IoT-Plattform}
        {}
        {}
        {\begin{itemize}
\item End-to-End-Produktentwicklung geleitet — von Hardware-Design über Lieferantenkoordination bis Cloud-Plattform-Deployment.
\item Embedded-KI für Tiererkennung und Verhaltenserkennung mit Micropython und MQTT implementiert.
\item Plattformübergreifende Mobile-App mit Expo (React Native) für Haustierüberwachung und Türsteuerung entwickelt.
\item Gesamten System-Stack aufgebaut: Firmware, Cloud-Kommunikation, Monitoring-Dashboard und mobile Benutzeroberfläche.
\item Domänenübergreifende Koordination zwischen Hardware-Prototyping, Softwareentwicklung und Nutzertests gesteuert.
\end{itemize}
\textbf{Schlüsselwörter}: IoT, Embedded-KI, Produktentwicklung, Hardware-Software-Integration}
\end{document}